% ==============================================================================
% DEFINIZIONE AMBIENTI TCOLORBOX & TEOREMI - FISICA GENERALE 1
% ==============================================================================

% --- Tavolozza Colori Istituzionali ed Eleganti ---
\definecolor{navyblue}{RGB}{15, 43, 92}
\definecolor{forestgreen}{RGB}{30, 92, 43}
\definecolor{purpleaccent}{RGB}{90, 32, 112}
\definecolor{crimsonred}{RGB}{139, 0, 0}
\definecolor{warmamber}{RGB}{184, 107, 0}
\definecolor{charcoalgray}{RGB}{43, 43, 43}
\definecolor{tealblue}{RGB}{0, 102, 119}
\definecolor{softgold}{RGB}{197, 155, 39}
\definecolor{darkcyan}{RGB}{0, 139, 139}

% Contatore globale per teoremi, definizioni, principi condiviso per capitolo
\newcounter{physcounter}[chapter]
\renewcommand{\thephyscounter}{\thechapter.\arabic{physcounter}}
\crefname{physcounter}{risultato}{risultati}
\Crefname{physcounter}{Risultato}{Risultati}

% --- 1. DEFINIZIONE ---
\newtcbtheorem[use counter=physcounter, number within=chapter]{definizione}{Definizione}{
    enhanced,
    breakable,
    colback=blue!3!white,
    colframe=navyblue,
    coltitle=white,
    fonttitle=\bfseries\sffamily,
    fontupper=\itshape,
    attach boxed title to top left={yshift=-2mm, xshift=4mm},
    boxed title style={
        sharp corners,
        rounded corners=north,
        colback=navyblue,
        boxrule=0pt
    },
    separator sign={:~},
    description font=\normalfont\bfseries,
    sharp corners,
    rounded corners=south,
    boxrule=1.2pt,
    top=4mm,
    bottom=3mm,
    left=3mm,
    right=3mm
}{def}

% --- 2. TEOREMA ---
\newtcbtheorem[use counter=physcounter, number within=chapter]{teorema}{Teorema}{
    enhanced,
    breakable,
    colback=forestgreen!3!white,
    colframe=forestgreen,
    coltitle=white,
    fonttitle=\bfseries\sffamily,
    fontupper=\itshape,
    attach boxed title to top left={yshift=-2mm, xshift=4mm},
    boxed title style={
        sharp corners,
        rounded corners=north,
        colback=forestgreen,
        boxrule=0pt
    },
    separator sign={:~},
    description font=\normalfont\bfseries,
    sharp corners,
    rounded corners=south,
    boxrule=1.2pt,
    top=4mm,
    bottom=3mm,
    left=3mm,
    right=3mm
}{teo}

% --- 3. PRINCIPIO / LEGGE FONDAMENTALE ---
\newtcbtheorem[use counter=physcounter, number within=chapter]{legge}{Principio / Legge}{
    enhanced,
    breakable,
    colback=purpleaccent!3!white,
    colframe=purpleaccent,
    coltitle=white,
    fonttitle=\bfseries\sffamily,
    fontupper=\itshape,
    attach boxed title to top left={yshift=-2mm, xshift=4mm},
    boxed title style={
        sharp corners,
        rounded corners=north,
        colback=purpleaccent,
        boxrule=0pt
    },
    separator sign={:~},
    description font=\normalfont\bfseries,
    sharp corners,
    rounded corners=south,
    boxrule=1.2pt,
    top=4mm,
    bottom=3mm,
    left=3mm,
    right=3mm
}{legge}

% --- 4. PROPOSIZIONE / COROLLARIO ---
\newtcbtheorem[use counter=physcounter, number within=chapter]{proposizione}{Proposizione}{
    enhanced,
    breakable,
    colback=tealblue!3!white,
    colframe=tealblue,
    coltitle=white,
    fonttitle=\bfseries\sffamily,
    fontupper=\itshape,
    attach boxed title to top left={yshift=-2mm, xshift=4mm},
    boxed title style={
        sharp corners,
        rounded corners=north,
        colback=tealblue,
        boxrule=0pt
    },
    separator sign={:~},
    description font=\normalfont\bfseries,
    sharp corners,
    rounded corners=south,
    boxrule=1.2pt,
    top=4mm,
    bottom=3mm,
    left=3mm,
    right=3mm
}{prop}

% --- 5. DIMOSTRAZIONE RIGOROSA ---
\newtcolorbox{dimostrazione}[1][]{
    enhanced,
    breakable,
    blanker,
    borderline west={2pt}{0pt}{forestgreen!80!black},
    left=4mm,
    top=2mm,
    bottom=2mm,
    before upper={\textbf{\textsf{\color{forestgreen!80!black}Dimostrazione#1:}}\par\vspace{1mm}},
    after upper={\par\hfill\textcolor{forestgreen!80!black}{$\blacksquare$}}
}

% --- 6. ESEMPIO APPLICATIVO ---
\newtcolorbox{esempio}[1][]{
    enhanced,
    breakable,
    colback=purpleaccent!3!white,
    colframe=purpleaccent!70!black,
    coltitle=white,
    fonttitle=\bfseries\sffamily,
    title={Esempio Fisico: #1},
    sharp corners,
    rounded corners=south,
    boxrule=1pt,
    top=3mm,
    bottom=3mm,
    left=3mm,
    right=3mm
}

% --- 7. ESERCIZIO GUIDATO SVOLTO ---
\newtcolorbox{esercizio}[1][]{
    enhanced,
    breakable,
    colback=warmamber!4!white,
    colframe=warmamber,
    coltitle=white,
    fonttitle=\bfseries\sffamily,
    title={Esercizio Guidato: #1},
    sharp corners,
    rounded corners=south,
    boxrule=1.1pt,
    top=3mm,
    bottom=3mm,
    left=3mm,
    right=3mm
}

% --- 8. CONSIGLI ESAME, ERRORI FREQUENTI E TRAPPOLE ---
\newtcolorbox{esame}[1][Consigli per la Prova d'Esame]{
    enhanced,
    breakable,
    colback=crimsonred!4!white,
    colframe=crimsonred,
    coltitle=white,
    fonttitle=\bfseries\sffamily,
    title={\textbf{Attenzione all'Esame:} #1},
    sharp corners,
    rounded corners=south,
    boxrule=1.2pt,
    top=3mm,
    bottom=3mm,
    left=3mm,
    right=3mm
}

% --- 9. METODO RISOLUTIVO / ALGORITMO OPERATIVO ---
\newtcolorbox{metodo}[1][]{
    enhanced,
    breakable,
    colback=warmamber!5!white,
    colframe=warmamber!90!black,
    coltitle=white,
    fonttitle=\bfseries\sffamily,
    title={Metodo Risolutivo Step-by-Step: #1},
    sharp corners,
    rounded corners=south,
    boxrule=1.2pt,
    top=3mm,
    bottom=3mm,
    left=3mm,
    right=3mm
}

% --- 10. OSSERVAZIONE / APPROFONDIMENTO ---
\newtcolorbox{osservazione}[1][Osservazione]{
    enhanced,
    breakable,
    blanker,
    borderline west={2.5pt}{0pt}{navyblue!70!black},
    left=4mm,
    top=2mm,
    bottom=2mm,
    before upper={\textbf{\textsf{\color{navyblue!70!black}#1:}}\par\vspace{1mm}}
}

% --- 11. ESPERIMENTO / FENOMENOLOGIA ---
\newtcolorbox{esperimento}[1][Esperimento / Fenomeno]{
    enhanced,
    breakable,
    colback=darkcyan!4!white,
    colframe=darkcyan,
    coltitle=white,
    fonttitle=\bfseries\sffamily,
    title={Fenomenologia Sperimentale: #1},
    sharp corners,
    rounded corners=south,
    boxrule=1pt,
    top=3mm,
    bottom=3mm,
    left=3mm,
    right=3mm
}
